\section{\textit{API}}

Com o objetivo de ampliar os casos de uso da ferramenta desenvolvida, a integração entre as camadas de \textit{frontend} e \textit{backend} é realizada por meio de uma \textit{API} pública, permitindo que sistemas externos também possam consumir as funcionalidades disponibilizadas. Essa abordagem favorece a interoperabilidade do \textit{software} e possibilita sua aplicação em diferentes contextos computacionais, como simulações, automações de processos e integração com outras plataformas utilizadas na engenharia química.

Para o consumo da \textit{API}, a documentação encontra-se disponível no \textit{endpoint} \textit{/docs} do \textit{backend}. Essa documentação foi desenvolvida utilizando o \textit{framework} \textit{FastAPI}, que possibilita a geração automática e padronizada da interface de documentação por meio da especificação \textit{OpenAPI}, facilitando a compreensão, o teste e a integração dos \textit{endpoints} disponibilizados.

No \textit{frontend}, o consumo dessa mesma \textit{API} é feito por um cliente tipado em \textit{TypeScript}, centralizado em uma camada única de serviços e exposto à interface por meio de chamadas ao caminho \textit{/api}, preservando a mesma origem no ambiente de desenvolvimento e no ambiente publicado. Dessa forma, a interface desenvolvida com \textit{React} e \textit{Vite} mantém a separação entre apresentação e dados, enquanto os módulos de tubulações, escoamento, dimensionamento, bombas, propriedades, reatores e balanço de massa acessam os mesmos \textit{endpoints} sem duplicar regras de negócio.

Um aspecto importante do estado atual do sistema é que a \textit{API} não entrega apenas resultados escalares. Em vários módulos, o \textit{backend} também monta estruturas de dados prontas para visualização, como modelos de gráfico com eixos, séries, marcadores e metadados de unidades. Esse padrão é empregado, por exemplo, no diagrama de \textit{Moody}, nas curvas de perda de carga e altura manométrica, no indicador de margem de \textit{NPSH}, no diagrama de \textit{Levenspiel}, no gráfico de \textit{Arrhenius}, nos perfis espaciais e de reciclo em \textit{PFR}, no diagrama ternário, no diagrama de \textit{McCabe-Thiele}, no envelope de fase, na curva de pressão de vapor e no gráfico nativo do balanço de massa.

Esse desenho arquitetural desloca para o \textit{backend} a responsabilidade por cálculos, amostragem numérica, construção de séries e definição dos pontos operacionais que precisam ser destacados. Ao \textit{frontend} cabe interpretar tais contratos tipados e renderizá-los em componentes \textit{SVG} reutilizáveis, preservando coerência visual entre páginas distintas. O mesmo princípio é adotado nos \textit{endpoints} de exemplo: a \textit{API} fornece entradas de referência coerentes com os catálogos disponíveis, e a interface apenas pré-preenche os formulários e dispara o fluxo normal de cálculo, sem manter resultados estáticos desacoplados do motor computacional.
